\section{Perspectives d'amélioration}
\begin{frame}[t]
    \frametitle{Perspectives d'amélioration à l'avenir}

    \small
    \textbf{Représentation au format SMILES.} 
    \begin{itemize}
        \item Utiliser une représentation chimique explicite peut enrichir la modélisation des structures.
        \item Meilleure précision pour la représentation 3D par la suite.
        \item Comparaison comparer peptide-only, SMILES-only et une approche hybride.
        \item Impact sur la génération/scoring.
    \end{itemize} 
    \vspace{0.45cm}
    \textbf{Utilisation d'un LLM pour l'input.}
    \begin{itemize}
        \item Langage naturel $\rightarrow$ contraintes structurées (propriétés cibles, seuils, priorités).
        \item Facilite l'utilisation du pipeline par des profils non techniques.
        \item Facilité de compréhension pour des contraintes biologiques.
    \end{itemize}
    \vspace{0.45cm}
    \textbf{Modèle spécialisé en peptides.}
    \begin{itemize}
        \item Embeddings adaptés au domaine $\rightarrow$ mieux capturer les motifs biochimiques pertinents.
        \item Améliorer la qualité des candidats générés.
        \item Résultats plus fiables car basés sur des motifs pertinents.
    \end{itemize}
\end{frame}